\documentclass[11pt]{article}
\usepackage{fontspec}
\usepackage[french]{babel}
\usepackage{amsmath}
\usepackage{amssymb}
\usepackage{amsthm}
\usepackage{mathrsfs}
\usepackage{tikz-cd}
\usepackage{enumitem}
\usetikzlibrary{babel}

\newcommand{\printedpage}[1]{\marginpar{\footnotesize p.~#1}}
\newcommand{\msp}[1]{\marginpar{\footnotesize [#1]}}
\newcommand{\Ob}{\mathsf{Ob}}
\newcommand{\id}{\mathsf{id}}
\newcommand{\DD}{\mathbb{D}}

\begin{document}

\printedpage{29}
% head: A. Grothendieck : Dérivateurs, Ch. I, Édition des Universités de Montpellier II et Paris VII
% foot: Version 8 juillet 2001

\noindent\msp{page 51}
changement de base pour
\[
\begin{tikzcd}
Z' \arrow[r, "p"] \arrow[d, "f'"'] & Z \arrow[d, "f"] \\
S \arrow[r, "\id"] & S
\end{tikzcd}
\qquad .
\]

\noindent\textsc{Remarque} : On retrouve le fait que $i^*i_*(F_Y) \overset{\sim}{\to} F_Y$ est \texttt{[un]} isomorphisme, pour $i$ une immersion fermée. De plus, avec les notations de \textbf{Der 4} (g), $j^*(i_*(F_Y))$ se calcule par le théorème de changement de base (corollaire 2 \texttt{[plutôt 4]} ci-dessus), qui s'applique dans toute situation commutative
\[
\begin{tikzcd}[row sep=large, column sep=large]
X' \arrow[r, "q"] \arrow[d, "f'"'] \arrow[dr, phantom, "\text{\scriptsize cart}"] & X \arrow[d, "f"] \\
Y' \arrow[r, "p"] & Y
\end{tikzcd}
\]
quand $p$ est une immersion ouverte (et $\gamma = \id$) (alors $X'/y' \to X/p(y')$ est même un isomorphisme!), et on trouve, revenant aux notations de \textbf{Der 4} (g),
\[
\begin{tikzcd}
\emptyset \arrow[r, "q"] \arrow[d, "f'"'] & Y \arrow[d, hook, "i"] \\
U \arrow[r, hook, "j"'] & X
\end{tikzcd}
\]
\textit{i.e.} $j^*(i_*(F_Y)) \simeq f'_*(\varphi)$, où $\varphi = q^*(F_Y)$ est l'unique élément de la catégorie (équivalente à une catégorie ponctuelle, par \textbf{Der 1}) $\DD(\emptyset)$. D'autre part, $f'_*$ transforme objet final en objet final, d'où $j^*i_*(F_Y)$ est objet final.

Inversement, soit $F$ dans $\DD(X)$, tel que $j^*(F) = e_{\DD(U)}$. Considérons
\[
F \xrightarrow{u} i_*i^*F \qquad\quad \text{(adjonction) ;}
\]
je dis que c'est un isomorphisme. En effet, par l'axiome \msp{page 52} de localisation \textbf{Der 2}, il suffit de prouver que $i^*(u)$ et $j^*(u)$ sont \texttt{[des]} isomorphismes. Pour $i^*(u)$ cela résulte du fait que $i^*$ soit foncteur de localisation, et pour $j^*$ du fait que $j^*(F)$ et $j^*(i_*i^*(F))$ sont des objets finaux dans $\DD(U)$.

\printedpage{30}
% head: A. Grothendieck : Dérivateurs, Ch. I, Édition des Universités de Montpellier II et Paris VII
% foot: Version 8 juillet 2001

\noindent\msp{page 53}

\subsection*{Morphismes homotopiquement propres, homotopiquement lisses.}

\noindent Soit $f : X \to Y$ \texttt{[un]} morphisme de catégories, considérons $y_0 \in \Ob\, Y$, et le foncteur canonique
\[
X_{y_0} \xrightarrow{\varphi} X/y_0 \qquad\qquad\qquad \varphi(x) = (x, \id_{f(x)}) \tag*{(1)}
\]
($X_{y_0}$ fibre de $X$ en $y_0$). Soit $(x_0, u_0) \in \Ob\, X/y_0$ donc $x_0 \in \Ob\, X$, $u_0 : f(x_0) \to y_0$. Considérons la catégorie coinduite par $\varphi$ sur $(x_0, u_0)$
\[
(x_0, u_0)\backslash(X_{y_0}) \quad ,
\]
pour tester si $\varphi$ est une hot-équivalence costricte \textit{i.e.} homologiquement \texttt{[plutôt cohomologiquement]} cofinal. On va l'interpréter via le foncteur induit par $f$
\[
x_0\backslash X \xrightarrow{\psi} f(x_0)\backslash Y
\]
allant de la catégorie des spécialisations de $x_0$ dans $X$ vers celle des spécialisations de $f(x_0)$ dans $Y$. On trouve alors un isomorphisme canonique de la catégorie coinduite $(x_0, u_0)\backslash X_{y_0}$ par $X_{y_0}$ au dessus de l'objet $(x_0, u_0)$ de $X/y_0$ et de la catégorie fibre $(x_0\backslash X)_{u_0}$ de $x_0\backslash X$ au dessus de $f(x_0)\backslash Y$ en $u_0$. Donc

\noindent\textbf{Proposition.} \textit{Pour que les foncteurs $\varphi : X_{y_0} \to X/y_0$ ($y_0 \in \Ob\, Y$) soient des hot-équivalences costrictes, il faut et il suffit que les fibres des foncteurs $x_0\backslash X \to f(x_0)\backslash Y$ induits par $f$ soient contractiles.}

\noindent\msp{page 54}
On dit alors que $f$ est \textit{homotopiquement propre.} La notion duale s'exprime, soit par la condition que les foncteurs induits par $f$
\[
X_{y_0} \longrightarrow y_0\backslash Y
\]
soient des hot-équivalences strictes \textit{i.e.} qu'ils soient cohomologiquement \texttt{[plutôt homologiquement]} cofinaux, soit que les fibres des morphismes induits sur les ``catégories localisées'' (catégories des générisations)
\[
X/x_0 \longrightarrow Y/f(x_0)
\]
soient contractiles. (N.B. On utilise ici qu'une catégorie $Z$ est contractile si et seulement si $Z^\circ$ l'est.) On dit alors que $f$ est \textit{homotopiquement lisse.}

\noindent\textbf{Proposition.} \textit{La notion de foncteur propre (resp. lisse) est stable par changement de base, par composition. Les isomorphismes, et plus généralement, les immersions fermées sont propres. L'image d'un morphisme propre est une sous-catégorie fermée.}\footnote[23]{Attention, une équivalence de catégories n'est pas propre! Si $X' \hookrightarrow X$ \texttt{[est]} une inclusion d'une sous-catégorie de $X$, elle n'est propre que si c'est une immersion fermée.}

\noindent N.B. Le foncteur $X \to e$ est toujours propre et lisse, car les catégories $x_0\backslash X$ et $X/x_0$ sont contractiles (ayant objet initial ou objet final). Donc tout foncteur de projection $X \times Y \xrightarrow{pr_2} Y$ est à la fois propre et lisse.

\printedpage{31}
% head: A. Grothendieck : Dérivateurs, Ch. I, Édition des Universités de Montpellier II et Paris VII
% foot: Version 8 juillet 2001

\noindent\msp{page 55}
\textsc{Démonstration}. Stabilité par changement de base :
\[
\begin{tikzcd}
X \arrow[d, "f"'] \arrow[dr, phantom, "\text{\scriptsize cart}"] & X' \arrow[l, "q"'] \arrow[d, "f'"] \\
Y & Y' \arrow[l, "p"]
\end{tikzcd}
\qquad ,
\]
d'où pour $x' \in \Ob\, X'$
\[
\begin{tikzcd}[row sep=large]
X/x \arrow[d] & X'/x' \arrow[l] \arrow[d] \\
Y/y & Y'/y' \arrow[l]
\end{tikzcd}
\qquad ,
\]
où $x$, $y'$, $y$ sont les images de $x'$ dans $X$, $Y'$, $Y$. On vérifie que ce diagramme est \textit{cartésien}, donc les fibres de $X'/x' \longrightarrow Y'/y'$ sont isomorphes à des fibres de $X/x \longrightarrow Y/y$, donc contractiles si celles-ci le sont.

Dans un cas $X \xrightarrow{f} Y$ on vérifie que pour un foncteur bilocalisé de $f$
\[
X/x \xrightarrow{\varphi=f/x} Y/y \quad ,
\]
tout ``bilocalisé''
\[
(X/x)/\xi \longrightarrow (Y/y)/\eta \qquad\qquad (\xi : x' \to x\,, \quad \eta : y' \to y)
\]
de $\varphi$ s'identifie encore à un bilocalisé de $f$
\[
X/x' \longrightarrow Y/y' \quad .
\]
Donc si $f$ est lisse, les bilocalisés $f/x$ ($x \in \Ob\, X$) le sont aussi. Notons d'autre part.

\noindent\textbf{Proposition.}\footnote[24]{\textit{cf.} amélioration dans cor. 2, page suivante.} \textit{Tout morphisme propre (ou dualement lisse) à fibres contractiles est une hot-équivalence.}

\noindent\msp{page 56}
Cela provient du fait que pour un tel foncteur propre, si $F$ est un préfaisceau, ou complexe de préfaisceaux abéliens sur $X$, la formation des $Rf_*(F)$ commute au passage aux fibres, et de même pour $R^1f_*(F)$ si $F$ est un préfaisceau en groupes (pas même abéliens).

\noindent\textbf{Corollaire 1.} \textit{Si $f : X \to Y$ est cohomologiquement propre à fibres contractiles, alors $Y$ est contractile si et seulement si $X$ l'est.}

Grâce à ceci, on trouve que le composé de deux foncteurs propres (lisses) est propre (lisse).

\printedpage{32}
% head: A. Grothendieck : Dérivateurs, Ch. I, Édition des Universités de Montpellier II et Paris VII
% foot: Version 8 juillet 2001

\noindent N.B. Contrairement à ce qui a lieu en topologie pour la notion de propreté, si $gf$ est propre et $g$ propre, $f$ ne l'est pas nécessairement, \textit{i.e.} un $S$-morphisme de $S$-catégories $X$, $Y$ propres sur $S$ n'est pas nécessairement propre. Ainsi une section de $X/S$, pour $X$ propre sur $S$, n'est pas nécessairement propre, \textit{i.e.} n'est pas nécessairement une immersion fermée. C'est déjà faux si $S = e$ (auquel cas $f : X \to e$ est toujours propre).

\noindent\textbf{Corollaire 2.} \textit{Un morphisme propre ou lisse, à fibres contractiles est à la fois une hot-équivalence stricte et costricte.}

En effet, on applique le corollaire 1 à $X/y \to Y/y$ ou $y\backslash X \to y\backslash Y$.

\noindent\textbf{Théorème} (g) \textit{Soit une situation de changement de base dans} \textsf{Dia}
\[
\begin{tikzcd}[row sep=large, column sep=huge]
X \arrow[d, "f"'] \arrow[dr, phantom, "\text{cart}"] & X' \arrow[l, "q"'] \arrow[d, "f'"] \\
Y & Y' \arrow[l, "p"]
\end{tikzcd}
\qquad , \qquad\qquad p^*f_*(F) \xrightarrow{c} f'_*q^*(F)
\]
\textit{($F \in \DD(X)$), $\DD$ dérivateur exact à gauche. Si $f$ est (homotopiquement) propre ou $p$ homotopiquement lisse, alors $c$ est un isomorphisme.}

\noindent\msp{page 57}
La question c'est si pour tout $y' \in Y'$, posant $y = p(y')$ et considérant le foncteur induit par $q$
\[
X/y \longleftarrow X'/y'
\]
le morphisme correspondant
\[
\Gamma_{\DD}(X/y, F) \longrightarrow \Gamma_{\DD}(X'/y', q^*F)
\]
est \texttt{[un]} isomorphisme. Dans le cas $f$ propre, on considère
\[
\begin{tikzcd}[column sep=large]
X_y \arrow[d, "\sim\ \ \text{costricte}"'] & X_{y'} \arrow[l, "\sim"'] \arrow[d, "\sim\ \ \text{costricte}"] \\
X/y & X'/y' \arrow[l]
\end{tikzcd}
\]
d'où la conclusion, puisque par \textbf{Der 5} (g), dans
\[
\begin{tikzcd}[row sep=large, column sep=large]
\Gamma(X_y) \arrow[r] & \Gamma(X_{y'}) \\
\Gamma(X/y) \arrow[u] \arrow[r] & \Gamma(X'/y') \arrow[u]
\end{tikzcd}
\]
les flèches verticales et l'horizontale supérieure sont \texttt{[des]} isomorphismes, donc aussi l'horizontale inférieure.

\printedpage{33}
% head: A. Grothendieck : Dérivateurs, Ch. I, Édition des Universités de Montpellier II et Paris VII
% foot: Version 8 juillet 2001

Dans le cas $p$ lisse, considérons du coté $Y$, $Y'$ le carré de changement de base % [sic] coté -> côté, missing circumflex (cf. "Du côté" below)
\[
\begin{tikzcd}[row sep=large, column sep=large]
Y & Y' \arrow[l] \\
Y/y \arrow[u] & Y'/y' \arrow[l] \arrow[u]
\end{tikzcd}
\]

\noindent\msp{page 58}
d'où le diagramme contigu
\[
\begin{tikzcd}[column sep=small]
X \arrow[dd] & & X' \arrow[ll] \arrow[dd] & \\
& X/y \arrow[ul] & & X'/y' \arrow[ll, crossing over, "\kappa"'] \arrow[ul] \arrow[dd] \\
Y & & Y' \arrow[ll] & \\
& Y/y \arrow[ul] \arrow[from=uu, crossing over] & & Y'/y' \arrow[ll, "\pi"] \arrow[ul]
\end{tikzcd}
\qquad ,
\]
où les carrés latéraux (verticaux) sont cartésiens. On note que $\pi$ est lisse et à fibres contractiles (car $p$ est lisse) donc $\kappa$ aussi. Donc $\kappa$ est une hot-équivalence costricte, donc $\Gamma(X/y) \overset{\sim}{\longrightarrow} \Gamma(X'/y')$, on gagne encore.

\noindent N.B. Sauf erreur, on a une réciproque du théorème, \texttt{[à]} savoir : Si $f : X \to Y$ est tel que le théorème de changement de base pour $f_*(F)$ est vrai pour \textit{tout} dérivateur (il suffit sans doute de \textsf{HOTAB} et \textsf{HOT}) et pour tout $Y' \to Y$, alors $f$ est homotopiquement propre. D'autre part, si $g : Y' \to Y$ est tel que le théorème de changement de base pour $g$ est vrai pout tout $X \to Y$ et tout $F \in \DD(X)$ ($\DD$ dérivateur exact à droite quelconque), alors $g$ est homotopiquement lisse. % [sic] pout tout -> pour tout, misprint

\noindent Du côté homologique les foncteurs lisses $f$ sont exactement ceux pour lesquels $f_!$ commute à tout changement de base. Les foncteurs propres $g$ sont exactement ceux pour lesquels le changement de base par $g$ commute à $f_!$, pour tout $f : X \to Y$ et tout $\DD$ dérivateur exact à gauche.

\noindent\msp{page 59}
Ainsi, les foncteurs propres $f : X \to Y$ se caractérisent de \textit{quatre} façons remarquables
\begin{itemize}[label={}, leftmargin=0pt, itemindent=2\parindent, labelsep=0.3em]
\item[\textit{a})] $f_*$ commute à tout changement de base ($\forall$ dérivateur $\DD$ exact à gauche);
\item[\textit{b})] le changement de base $f^*$ commute à tout $g_!$ ($g : Y' \to Y$, $\DD$ dérivateur exact à droite);
\item[\textit{c})] $\forall y \in \Ob\, Y$, $X_y \to X/y$ est une hot-équivalence costricte \textit{i.e.} est un foncteur homologiquement \texttt{[plutôt cohomologiquement]} cofinal;
\item[\textit{d})] $\forall x \in \Ob\, X$ le foncteur induit sur les colocalisés
\[
x\backslash X \longrightarrow f(x)\backslash Y
\]
\end{itemize}

\end{document}
